\documentclass[11pt,a4paper]{article}
\usepackage[T1]{fontenc}
\usepackage[utf8]{inputenc}
\usepackage{amsmath,amssymb,amsthm}
\usepackage[margin=1in]{geometry}
\usepackage[colorlinks=true,linkcolor=blue,urlcolor=blue]{hyperref}
\theoremstyle{plain}
\newtheorem{theorem}{Theorem}
\newtheorem{lemma}[theorem]{Lemma}
\newtheorem{proposition}[theorem]{Proposition}
\newtheorem{corollary}[theorem]{Corollary}
\theoremstyle{definition}
\newtheorem{remark}[theorem]{Remark}

\title{Recovering the unwritten recurrences of the table OEIS A236809}
\author{Adrian Perez Fontelles\\ \small Independent researcher}
\date{2 September 2026}
\begin{document}
\maketitle

\begin{abstract}
OEIS A236809 is a two-dimensional table $T(n,k)$ whose empirical recurrences are, for several of
its lines, given only as an ORDER --- ``k=2: [order 10]'' and the like --- with the
coefficients in a linked file rather than in the entry. Those conjectures can still be settled,
and the recurrences recovered. A line of the table is a fixed-width or fixed-height array
count, hence a walk count on finitely many vertices, hence satisfies some monic recurrence of
order at most that number; Berlekamp--Massey on twice as many exact terms returns the MINIMAL
one. Where its order equals the order the entry states --- and where the same holds after the
first few terms are dropped --- any recurrence of that order the line satisfies has a
characteristic polynomial that is a multiple of the minimal one and of equal degree, hence is
that polynomial. This note settles column 2, column 3, column 4, column 5.
\end{abstract}

\noindent\small 2020 Mathematics Subject Classification. 05A15, 11B37, 15A18.\normalsize

\section{The table and the unwritten conjectures}

OEIS A236809 is ``T(n,k)=Number of (n+1)X(k+1) 0..2 arrays with the maximum plus the lower median minus the upper median minus the minimum of every 2X2 subblock equal.''. It is read by antidiagonals and begins
\[
81, 277, 277, 1033, 1435, 1033, 4183, 8825,\ \dots
\]
The entry carries, among its empirical claims:
\begin{quote}\small
k=2: [order 10]\\
k=3: [order 11]\\
k=4: [order 13]\\
k=5: [order 13]
\end{quote}
As of the ``Last modified'' line on the live entry (Jul 23 2025 09:34:47, revision 6) these are still
recorded as empirical, and the recurrences themselves are not in the entry text.

\section{Why a line of the table is a walk count}

Fix the index that the line fixes. The entry defines $T(n,k)$ as a number of arrays of a shape
determined by $n$ and $k$, subject to a condition local to a bounded window of consecutive
lines. Substituting the fixed index into the entry's own wording turns the definition into an
ordinary one-parameter array count, and for such a count the admissible windows are the
vertices of a finite digraph and an array is a walk. So the line is a walk count on $S$
vertices, and by the Cayley--Hamilton theorem it satisfies a monic integer recurrence of order
at most $S$.

\section{Recovering the recurrences}

\begin{lemma}\label{lem:bm}
A sequence known to satisfy some linear recurrence of order at most $S$ is determined, with its
minimal recurrence, by its first $2S$ terms; Berlekamp--Massey applied to them returns that
minimal recurrence exactly.
\end{lemma}

\subsection*{Column $k=2$}
Substituting the fixed index into the entry's wording gives an ordinary one-parameter array
count, a walk on $S=81$ vertices. Berlekamp--Massey on $2S=162$ exact terms
returns a minimal recurrence of order $10$ --- the order the entry states --- with
integer coefficients
\begin{center}\small\begin{tabular}{cccccccc}
$c_{1}$ & $12$ & $c_{4}$ & $2434$ & $c_{7}$ & $24192$ & $c_{10}$ & $23328$ \\
$c_{2}$ & $15$ & $c_{5}$ & $864$ & $c_{8}$ & $14472$ &  &  \\
$c_{3}$ & $-672$ & $c_{6}$ & $-17988$ & $c_{9}$ & $-46656$ &  &  \\
\end{tabular}\end{center}
so that $a(m)=\sum_{i=1}^{10}c_i\,a(m-i)$ for every $m$. Removing the first
$10+4$ terms and repeating gives order $10$ again, which is what the
identification below needs.

\subsection*{Column $k=3$}
Substituting the fixed index into the entry's wording gives an ordinary one-parameter array
count, a walk on $S=235$ vertices. Berlekamp--Massey on $2S=470$ exact terms
returns a minimal recurrence of order $11$ --- the order the entry states --- with
integer coefficients
\begin{center}\small\begin{tabular}{cccccccc}
$c_{1}$ & $24$ & $c_{4}$ & $26568$ & $c_{7}$ & $644112$ & $c_{10}$ & $1612416$ \\
$c_{2}$ & $-55$ & $c_{5}$ & $-71616$ & $c_{8}$ & $-762032$ & $c_{11}$ & $-783360$ \\
$c_{3}$ & $-2772$ & $c_{6}$ & $-74996$ & $c_{9}$ & $-588288$ &  &  \\
\end{tabular}\end{center}
so that $a(m)=\sum_{i=1}^{11}c_i\,a(m-i)$ for every $m$. Removing the first
$11+4$ terms and repeating gives order $11$ again, which is what the
identification below needs.

\subsection*{Column $k=4$}
Substituting the fixed index into the entry's wording gives an ordinary one-parameter array
count, a walk on $S=675$ vertices. Berlekamp--Massey on $2S=1350$ exact terms
returns a minimal recurrence of order $13$ --- the order the entry states --- with
integer coefficients
\begin{center}\small\begin{tabular}{cccccccc}
$c_{1}$ & $36$ & $c_{5}$ & $2245440$ & $c_{9}$ & $-974576448$ & $c_{13}$ & $1330255872$ \\
$c_{2}$ & $213$ & $c_{6}$ & $-34534356$ & $c_{10}$ & $1350053568$ &  &  \\
$c_{3}$ & $-20688$ & $c_{7}$ & $121902336$ & $c_{11}$ & $778781952$ &  &  \\
$c_{4}$ & $159340$ & $c_{8}$ & $49292928$ & $c_{12}$ & $-2623560192$ &  &  \\
\end{tabular}\end{center}
so that $a(m)=\sum_{i=1}^{13}c_i\,a(m-i)$ for every $m$. Removing the first
$13+4$ terms and repeating gives order $13$ again, which is what the
identification below needs.

\subsection*{Column $k=5$}
Substituting the fixed index into the entry's wording gives an ordinary one-parameter array
count, a walk on $S=1923$ vertices. Berlekamp--Massey on $2S=3846$ exact terms
returns a minimal recurrence of order $13$ --- the order the entry states --- with
integer coefficients
\begin{center}\small\begin{tabular}{cccccccc}
$c_{1}$ & $6$ & $c_{5}$ & $7631304$ & $c_{9}$ & $9534803520$ & $c_{13}$ & $-16293888000$ \\
$c_{2}$ & $2037$ & $c_{6}$ & $222376028$ & $c_{10}$ & $-17365107200$ &  &  \\
$c_{3}$ & $-12282$ & $c_{7}$ & $-1410692088$ & $c_{11}$ & $-5340748800$ &  &  \\
$c_{4}$ & $-1251404$ & $c_{8}$ & $774758880$ & $c_{12}$ & $29872128000$ &  &  \\
\end{tabular}\end{center}
so that $a(m)=\sum_{i=1}^{13}c_i\,a(m-i)$ for every $m$. Removing the first
$13+4$ terms and repeating gives order $13$ again, which is what the
identification below needs.

\begin{theorem}
For each line above, the displayed recurrence holds for every index, and it is the recurrence
the entry records.
\end{theorem}

\begin{proof}
It holds by Lemma~\ref{lem:bm}. For the identification, the entry asserts that the line
satisfies SOME recurrence of the stated order, possibly only from some index on. Let $q$ be its
characteristic polynomial and $q_{\min}$ the minimal polynomial of the tail. Then
$q_{\min}\mid q$, and the second Berlekamp--Massey run --- on the line with its first few
terms removed --- shows $\deg q_{\min}$ equals the stated order, which is $\deg q$. Both are
monic, so $q=q_{\min}$.
\end{proof}

\section{Verification}

Three checks.

First, each line's model was compared against the entry's own published values for that line,
taken from the antidiagonal DATA read in either orientation and from the ``Table starts''
block, and agrees at every available term. This is the only point at which reading the entry's
English enters, and a misreading gives different counts.

Second, each recovered recurrence was evaluated directly on the model's values, with no
Berlekamp--Massey involved, and holds at every index tested.

Third, each order was computed twice by independent routes --- modulo a $61$-bit prime, where
every intermediate is a machine word, and again in exact rational arithmetic --- and once more
on the line with its first few terms dropped, which is what the identification requires. All
agree.

\begin{thebibliography}{9}
\bibitem{oeis} The OEIS Foundation, \emph{The On-Line Encyclopedia of Integer
Sequences}, \url{https://oeis.org}, 2026. Sequence A236809.
\bibitem{massey} J.~L.~Massey, Shift-register synthesis and BCH decoding, \emph{IEEE
Trans. Inform. Theory} \textbf{15} (1969), 122--127.
\bibitem{stanley} R.~P.~Stanley, \emph{Enumerative Combinatorics}, Volume 1, 2nd ed.,
Cambridge University Press, 2011. (Transfer-matrix method, Section 4.7.)
\bibitem{horn} R.~A.~Horn and C.~R.~Johnson, \emph{Matrix Analysis}, 2nd ed., Cambridge
University Press, 2013.
\end{thebibliography}

\end{document}
